% GENERATED FILE. Edit canonical lessons, then run npm run generate:books.
% canonical-source-hash: fnv1a64:207d4bc469dd167c
% canonical-lessons: RU-C477-toska, RU-C477-toskovat, RU-C477-odinochestvo, RU-C477-odinokiy, RU-C477-sozhalet

\chapter{Longing and Loneliness}
\label{ch:ru-longing-and-loneliness}
\hlchaptermodality{477}

\begin{chapteropening}
\textbf{By the end of this chapter:} \emph{I can talk about missing people and places.}

Complete the last lesson of chapter 477: I can talk about missing people and places.
\end{chapteropening}

\section[\texorpdfstring{\ru{тоска} (toská)}{тоска (toská)}]{\ru{тоска} (toská) — longing; melancholy}

\label{lesson:RU-C477-toska}



\begin{quote}
Before the new one: say the Russian for a coward, then the Russian for courage.
\end{quote}

\subsection*{You'll want to know: \ru{тоска}}
\textbf{\ru{тоска}} (\emph{toská}) — \textquotedblleft{}longing; melancholy\textquotedblright{}.

\begin{quote}
\ru{Она} \ru{чувствует} \ru{тоску} \ru{по} \ru{родине}.
\end{quote}

\begin{quote}
\emph{Oná chúvstvuyet toskú po ródine.}
\end{quote}

\begin{quote}
She feels homesick.
\end{quote}

Feminine. \emph{\ru{Тоска} \ru{по}} + dative — longing for.

\begin{cousinweb}
\emph{\ru{тоска}} is inherited from Proto-Slavic \emph{t\ru{ъ}ska} (EtymDB), which EtymDB glosses 'emptiness'.
\end{cousinweb}

\begin{grammarlens}[title={What you've built}]
One more word for telling, explaining and describing at B1.
\end{grammarlens}

\subsection*{Your turn}
\begin{itemize}
\raggedright
  \item \emph{Say it:} \emph{toská}
  \item \emph{Say it:} \emph{toská}, once more
  \item \emph{Say it:} a sentence of your own with \emph{toská}
  \item \emph{Recall:} say the Russian for a coward, then the Russian for courage, then say \emph{toská} again
\end{itemize}

\begin{tcolorbox}[breakable,colback=teal!4,colframe=teal!35!black,title={Before you move on}]
What does \ru{тоска} mean? (\textquotedblleft{}Longing; melancholy\textquotedblright{}.) Say it once more.
\end{tcolorbox}
\section[\texorpdfstring{\ru{тосковать} (toskovát)}{тосковать (toskovát)}]{\ru{тосковать} (toskovát) — to pine, to miss}

\label{lesson:RU-C477-toskovat}



\begin{quote}
Before the new one: say the Russian for courage, then the Russian for longing.
\end{quote}

\subsection*{You'll want to know: \ru{тосковать}}
\textbf{\ru{тосковать}} (\emph{toskovát}) — \textquotedblleft{}to pine, to miss\textquotedblright{}.

\begin{quote}
\ru{Бабушка} \ru{тоскует} \ru{по} \ru{деревне}.
\end{quote}

\begin{quote}
\emph{Bábushka toskúyet po derévne.}
\end{quote}

\begin{quote}
Grandma misses the village.
\end{quote}

Imperfective; present \ru{тоскую}, \ru{тоскуешь}. \emph{\ru{Тосковать} \ru{по}} + dative.

\begin{cousinweb}
\emph{\ru{тосковать}} is built on \emph{\ru{тоска}} 'longing' with \emph{-\ru{овать}}; \emph{\ru{тоска}} is inherited from Proto-Slavic \emph{t\ru{ъ}ska} (EtymDB).
\end{cousinweb}

\begin{grammarlens}[title={What you've built}]
One more word for telling, explaining and describing at B1.
\end{grammarlens}

\subsection*{Your turn}
\begin{itemize}
\raggedright
  \item \emph{Say it:} \emph{toskovát}
  \item \emph{Say it:} \emph{toskovát}, once more
  \item \emph{Say it:} a sentence of your own with \emph{toskovát}
  \item \emph{Recall:} say the Russian for courage, then the Russian for longing, then say \emph{toskovát} again
\end{itemize}

\begin{tcolorbox}[breakable,colback=teal!4,colframe=teal!35!black,title={Before you move on}]
What does \ru{тосковать} mean? (\textquotedblleft{}To pine, to miss\textquotedblright{}.) Say it once more.
\end{tcolorbox}
\section[\texorpdfstring{\ru{одиночество} (odinóchestvo)}{одиночество (odinóchestvo)}]{\ru{одиночество} (odinóchestvo) — loneliness}

\label{lesson:RU-C477-odinochestvo}



\begin{quote}
Before the new one: say the Russian for longing, then the Russian for to pine.
\end{quote}

\subsection*{You'll want to know: \ru{одиночество}}
\textbf{\ru{одиночество}} (\emph{odinóchestvo}) — \textquotedblleft{}loneliness\textquotedblright{}.

\begin{quote}
\ru{После} \ru{переезда} \ru{он} \ru{долго} \ru{страдал} \ru{от} \ru{одиночества}.
\end{quote}

\begin{quote}
\emph{Pósle pereyézda on dólgo stradál ot odinóchestva.}
\end{quote}

\begin{quote}
After the move he suffered from loneliness for a long time.
\end{quote}

Neuter.

\begin{cousinweb}
\emph{\ru{одиночество}} is built on \emph{\ru{одинокий}} 'lonely', from \emph{\ru{один}} 'one', which is inherited from Proto-Slavic \emph{edin\ru{ъ}} (EtymDB).
\end{cousinweb}

\begin{grammarlens}[title={What you've built}]
One more word for telling, explaining and describing at B1.
\end{grammarlens}

\subsection*{Your turn}
\begin{itemize}
\raggedright
  \item \emph{Say it:} \emph{odinóchestvo}
  \item \emph{Say it:} \emph{odinóchestvo}, once more
  \item \emph{Say it:} a sentence of your own with \emph{odinóchestvo}
  \item \emph{Recall:} say the Russian for longing, then the Russian for to pine, then say \emph{odinóchestvo} again
\end{itemize}

\begin{tcolorbox}[breakable,colback=teal!4,colframe=teal!35!black,title={Before you move on}]
What does \ru{одиночество} mean? (\textquotedblleft{}Loneliness\textquotedblright{}.) Say it once more.
\end{tcolorbox}
\section[\texorpdfstring{\ru{одинокий} (odinókiy)}{одинокий (odinókiy)}]{\ru{одинокий} (odinókiy) — lonely}

\label{lesson:RU-C477-odinokiy}



\begin{quote}
Before the new one: say the Russian for to pine, then the Russian for loneliness.
\end{quote}

\subsection*{You'll want to know: \ru{одинокий}}
\textbf{\ru{одинокий}} (\emph{odinókiy}) — \textquotedblleft{}lonely\textquotedblright{}.

\begin{quote}
\ru{Наш} \ru{сосед} — \ru{одинокий} \ru{старик}.
\end{quote}

\begin{quote}
\emph{Nash soséd — odinókiy starík.}
\end{quote}

\begin{quote}
Our neighbour is a lonely old man.
\end{quote}

Also single, living alone. Short forms: \ru{одинок}, \ru{одинока}.

\begin{cousinweb}
\emph{\ru{одинокий}} is built on \emph{\ru{один}} 'one', which is inherited from Proto-Slavic \emph{edin\ru{ъ}}, from PIE \emph{óynos 'one' (EtymDB).}
\end{cousinweb}

\begin{grammarlens}[title={What you've built}]
One more word for telling, explaining and describing at B1.
\end{grammarlens}

\subsection*{Your turn}
\begin{itemize}
\raggedright
  \item \emph{Say it:} \emph{odinókiy}
  \item \emph{Say it:} \emph{odinókiy}, once more
  \item \emph{Say it:} a sentence of your own with \emph{odinókiy}
  \item \emph{Recall:} say the Russian for to pine, then the Russian for loneliness, then say \emph{odinókiy} again
\end{itemize}

\begin{tcolorbox}[breakable,colback=teal!4,colframe=teal!35!black,title={Before you move on}]
What does \ru{одинокий} mean? (\textquotedblleft{}Lonely\textquotedblright{}.) Say it once more.
\end{tcolorbox}
\section[\texorpdfstring{\ru{сожалеть} (sozhalét)}{сожалеть (sozhalét)}]{\ru{сожалеть} (sozhalét) — to regret}

\label{lesson:RU-C477-sozhalet}



\begin{quote}
Before the new one: say the Russian for loneliness, then the Russian for lonely.
\end{quote}

\subsection*{You'll want to know: \ru{сожалеть}}
\textbf{\ru{сожалеть}} (\emph{sozhalét}) — \textquotedblleft{}to regret\textquotedblright{}.

\begin{quote}
\ru{Я} \ru{сожалею}, \ru{что} \ru{не} \ru{поехал} \ru{с} \ru{вами}.
\end{quote}

\begin{quote}
\emph{Ya sozhaléyu, chto ne poyékhal s vámi.}
\end{quote}

\begin{quote}
I regret not going with you.
\end{quote}

Imperfective; present \ru{сожалею}. \emph{\ru{Сожалеть} \ru{о}} + prepositional. More formal than \emph{\ru{жалеть}}.

\begin{cousinweb}
\emph{\ru{сожалеть}} is built on \emph{\ru{со}-} and \emph{\ru{жалеть}} 'to pity', from \emph{\ru{жаль}} 'pity'; EtymDB lists these words without a source, and they are usually traced to Proto-Slavic \emph{žal\ru{ь}} 'grief'.
\end{cousinweb}

\begin{grammarlens}[title={What you've built}]
One more word for telling, explaining and describing at B1.
\end{grammarlens}

\subsection*{Your turn}
\begin{itemize}
\raggedright
  \item \emph{Say it:} \emph{sozhalét}
  \item \emph{Say it:} \emph{sozhalét}, once more
  \item \emph{Say it:} a sentence of your own with \emph{sozhalét}
  \item \emph{Recall:} say the Russian for loneliness, then the Russian for lonely, then say \emph{sozhalét} again
\end{itemize}

\begin{tcolorbox}[breakable,colback=teal!4,colframe=teal!35!black,title={Before you move on}]
What does \ru{сожалеть} mean? (\textquotedblleft{}To regret\textquotedblright{}.) Say it once more.
\end{tcolorbox}
